\bookchapter{Your workstation}\label{k:ws}

Everything in this book starts at a laptop: you write the code there, you test it against a simulated robot there, and you watch the real robot from there. This chapter sets up that laptop once, so that every later chapter can say ``open a terminal in the repository'' and it simply works. Milestone: both.

\begin{needbox}
\begin{itemize}[nosep]
  \item A laptop or desktop with \textbf{Ubuntu 26.04 LTS} (other recent Ubuntu versions work too; the commands are the same).
  \item For training policies (Chapters~\ref{k:walk} and~\ref{k:train}): an \textbf{NVIDIA graphics card}. Without one you can do every chapter except the training runs; use a shared training computer of the team for those.
  \item A GitHub account, and about $20$\,GB of free disk space.
  \item Background: \rtr{Bash Basics}, \rtr{Packages and Environments}, \rtr{Git and GitHub}.
\end{itemize}
\end{needbox}

\section*{The plan}
\begin{center}
\begin{tikzpicture}[>={Stealth[length=2.2mm]},bx/.style={draw,thick,rounded corners=3pt,minimum width=3.0cm,minimum height=0.9cm,font=\scriptsize,align=center}]
 \node[bx,fill=blue!8] (lap) at (0,0) {laptop\\ code, tests, robotcontrol};
 \node[bx,fill=green!10] (sim) at (4.6,0.9) {fake robot / MuJoCo\\ (Chapter~\ref{k:fake})};
 \node[bx,fill=orange!12] (k1) at (4.6,-0.9) {K1 onboard computer\\ runs \texttt{k1code}};
 \node[bx,fill=gray!12] (gpu) at (-4.6,0) {GPU computer\\ trains policies};
 \draw[<->,thick] (lap)--(sim); \draw[<->,thick] (lap)--node[below left,font=\scriptsize]{SSH, debug} (k1);
 \draw[->,thick] (gpu)--node[above,font=\scriptsize]{\texttt{.onnx}} (lap);
\end{tikzpicture}
\captionof{figure}{The computers of the book. Code is written on the laptop, tested against a fake robot, and deployed to the K1. Policies are trained on a computer with a GPU (it may be the same laptop).}
\end{center}

\section*{Steps}
\begin{stepbox}{Step 1: basic tools}
\begin{verbatim}
sudo apt update && sudo apt upgrade
sudo apt install git curl tmux htop build-essential
\end{verbatim}
\end{stepbox}

\begin{stepbox}{Step 2: the graphics driver (only with an NVIDIA card)}
\begin{verbatim}
sudo ubuntu-drivers install
sudo reboot
nvidia-smi
\end{verbatim}
\verb!nvidia-smi! prints a table with the name of the card, the driver version and the memory. If it prints \emph{command not found} or \emph{no devices}, the driver is not loaded.
\end{stepbox}

\begin{stepbox}{Step 3: uv, the Python manager}
\begin{verbatim}
curl -LsSf https://astral.sh/uv/install.sh | sh
uv --version
\end{verbatim}
\texttt{uv} installs Python versions, creates virtual environments and installs packages from a lock file. We use it everywhere, on the laptop and on the robot, so that both run \emph{the same} versions (\rtr{Packages and Environments}).
\end{stepbox}

\begin{stepbox}{Step 4: git and an SSH key for GitHub}
\begin{verbatim}
git config --global user.name  "Your Name"
git config --global user.email "you@example.org"
ssh-keygen -t ed25519
cat ~/.ssh/id_ed25519.pub   # paste into GitHub -> Settings -> SSH keys
ssh -T git@github.com
\end{verbatim}
\end{stepbox}

\begin{stepbox}{Step 5: the code of this book}
\begin{verbatim}
mkdir -p ~/k1 && cd ~/k1
git clone git@github.com:<your-team>/k1-in-practice-code.git
cd k1-in-practice-code
uv sync          # creates .venv and installs the locked packages
uv run pytest    # runs the tests; all must pass
\end{verbatim}
The address is the one of your team's copy (a \emph{fork}) of the code repository\guess. Chapter~\ref{k:repo} explains what is inside.
\end{stepbox}

\begin{stepbox}{Step 6: the simulator}
The physics simulator of the book is MuJoCo; training uses mjlab (\rtr{Your first mjlab environment}). The quickest test from the mjlab documentation is
\begin{verbatim}
uvx --from mjlab --refresh demo
\end{verbatim}
A window opens with a simulated robot. Close it. Then get the Booster tasks and the K1 model:
\begin{verbatim}
cd ~/k1
git clone https://github.com/BoosterRobotics/booster-mjlab.git
cd booster-mjlab && uv sync
\end{verbatim}
The repository address and the folder that holds the K1 model (an MJCF file with the $22$ joints) are not verified\guess; look for a file with \texttt{k1} in its name and \texttt{.xml} at the end: \verb!find . -iname "*k1*.xml"!.
\end{stepbox}

\begin{stepbox}{Step 7: look at the K1}
\begin{verbatim}
uv run python -m mujoco.viewer --mjcf=<path to the K1 xml>
\end{verbatim}
In the viewer, press the space bar to start and stop the simulation, double-click a body to select it, and use the panel on the right to move the joints with sliders.
\end{stepbox}

\begin{checkbox}
\begin{itemize}[nosep]
  \item \verb!uv run pytest! in the code repository ends with \emph{passed} and no \emph{failed}.
  \item The K1 appears in the MuJoCo viewer. Without a controller it collapses when you start the simulation: that is correct, nothing holds the joints yet. With the \emph{keyframe} ``home'' (or similar)\guess{} loaded and the simulation paused, it stands in its default pose.
  \item \verb!nvidia-smi! shows your card (only if you have one).
\end{itemize}
\end{checkbox}

\begin{troublebox}
\begin{description}[leftmargin=0.6cm,nosep]
  \item[\texttt{uv: command not found}] Open a new terminal; the installer adds \verb!~/.local/bin! to the \texttt{PATH} only for new shells.
  \item[The viewer does not open] On a computer without a screen (a server over SSH) there is no window. Use the laptop, or forward the display with \verb!ssh -X!.
  \item[\texttt{Permission denied (publickey)}] The public key is not in GitHub, or you used the HTTPS address. Run \verb!ssh -vT git@github.com! and read which key is offered.
  \item[Tests fail right after cloning] Check the Python version: \verb!uv run python --version! must match the one in \texttt{pyproject.toml}.
\end{description}
\end{troublebox}

\section*{ROS 2 on the laptop?}
You do \emph{not} need ROS~2 on the laptop for this book. ROS~2 runs on the robot and only delivers sensor data into our program (Chapter~\ref{k:ros}); the laptop talks to the robot through our own debug connection (Chapter~\ref{k:rc}). This is a deliberate choice: ROS~2 on two computers with different Ubuntu versions is a common source of trouble (different ROS distributions do not always talk to each other).

\begin{keyidea}[What to remember]
\begin{itemize}[nosep]
  \item One laptop with git, uv and (if possible) an NVIDIA driver; the same Python versions on laptop and robot thanks to the lock file.
  \item \verb!uv sync! then \verb!uv run pytest! is the first thing you do after every clone or pull.
  \item The simulator shows the K1 model; without a controller the model falls, and that is expected.
\end{itemize}
\end{keyidea}

\section*{Practice problems}
\begin{enumerate}[label=\textbf{\thechapter.\arabic*}]
  \item Why does the book insist on a \emph{lock file} (\texttt{uv.lock}) instead of a list of package names? Describe a bug that a lock file prevents.
  \item A training run needs $200$ million simulation steps. Your GPU simulates $4096$ robots at $50$ policy steps per second each in real time, and the simulation runs $10$ times faster than real time. How long does the training take?
  \item \verb!nvidia-smi! shows \emph{7800MiB / 8192MiB} in use while nothing is training. Name a likely cause and a command to find it.
  \item Write the two git commands that bring your copy of the code up to date and then show which tag you are on.
  \item Explain in two sentences why the laptop does not need ROS~2 in this book.
\end{enumerate}
